\documentclass[a4paper, 12pt]{article}

%%% Работа с русским языком
\usepackage{cmap}					% поиск в PDF
\usepackage{mathtext} 				% русские буквы в формулах
\usepackage[T2A]{fontenc}			% кодировка
\usepackage[utf8]{inputenc}			% кодировка исходного текста
\usepackage[russian]{babel}	% локализация и переносы

%%% Дополнительная работа с математикой
\usepackage{amsmath,amsfonts,amssymb,amsthm,mathtools} % AMS
\usepackage{icomma} % "Умная" запятая: $0,2$ --- число, $0, 2$ --- перечисление

%% Номера формул
%\mathtoolsset{showonlyrefs=true} % Показывать номера только у тех формул, на которые есть \eqref{} в тексте.

%% Шрифты
\usepackage{euscript}	 % Шрифт Евклид
\usepackage{mathrsfs} % Красивый матшрифт

%% Поля
\usepackage[left=2cm,right=2cm,top=2cm,bottom=2cm,bindingoffset=0cm]{geometry}

%% Русские списки
\usepackage{enumitem}
\makeatletter
\AddEnumerateCounter{\asbuk}{\russian@alph}{щ}
\makeatother

%%% Работа с картинками
\usepackage{graphicx}  % Для вставки рисунков
\graphicspath{{images/}{images2/}}  % папки с картинками
\setlength\fboxsep{3pt} % Отступ рамки \fbox{} от рисунка
\setlength\fboxrule{1pt} % Толщина линий рамки \fbox{}
\usepackage{wrapfig} % Обтекание рисунков и таблиц текстом

%%% Работа с таблицами
\usepackage{array,tabularx,tabulary,booktabs} % Дополнительная работа с таблицами
\usepackage{longtable}  % Длинные таблицы
\usepackage{multirow} % Слияние строк в таблице

%% Красная строка
\setlength{\parindent}{2em}

%% Интервалы
\linespread{1}
\usepackage{multirow}

%% TikZ
\usepackage{tikz}
\usetikzlibrary{graphs,graphs.standard}

%% Верхний колонтитул
\usepackage{fancyhdr}
\pagestyle{fancy}

%% Перенос знаков в формулах (по Львовскому)
\newcommand*{\hm}[1]{#1\nobreak\discretionary{}
	{\hbox{$\mathsurround=0pt #1$}}{}}

%% Мои дополнения
\usepackage{float} %Добавляет возможность работы с командой [H] которая улучшает расположение на странице
\usepackage{gensymb} %Красивые градусы
\usepackage{graphicx}               % Импорт изображений
\usepackage{caption} % Пакет для подписей к рисункам, в частности, для работы caption*

% подключаем hyperref (для ссылок внутри  pdf)
\usepackage[unicode, pdftex]{hyperref}

%%% Теоремы
\theoremstyle{plain}                    % Это стиль по умолчанию, его можно не переопределять.
\renewcommand\qedsymbol{$\blacksquare$} % переопределение символа завершения доказательства

\newtheorem{theorem}{Теорема}[section] % Теорема (счетчик по секиям)
\newtheorem{proposition}{Утверждение}[section] % Утверждение (счетчик по секиям)
\newtheorem{definition}{Определение}[section] % Определение (счетчик по секиям)
\newtheorem{corollary}{Следствие}[theorem] % Следстиве (счетчик по теоремам)
\newtheorem{problem}{Задача}[section] % Задача (счетчик по секиям)
\newtheorem*{remark}{Примечание} % Примечание (можно переопределить, как Замечание)
\newtheorem{lemma}{Лемма}[section] % Лемма (счетчик по секиям)

\begin{document}
	\newcommand{\HRule}{\rule{\linewidth}{0.7mm}} % Defines a new command for the horizontal lines, change thickness here
	
	\begin{center}
		\large\textbf{Московский Физико-Технический Институт}\\ % Name of your university/college
		\large\textbf{(государственный университет)}
	
		\vfill
		
		\Large Лабораторная работа по курсу общей физики № *labnum*\\[0.5cm] % Preambule of your document title
		
		
		\HRule
		\\[0.4cm]
		{ \huge \bfseries *name of your labwork*}% Title of your document
		\\[0.4cm] 
		\HRule
		\\[0.5cm]
		
		\ \\
	\textbf{\large Автор:} \\	
	\large *your name* *groupname*\\ % Your name and something more, your group num for example
		\vfill
		\hspace*{-0.8 cm}\includegraphics[width=100 pt]{frkt_logo}\\ % logo of your  company/university/college
		\large Долгопрудный, 2021 % location and year
	\end{center}

\newpage
\setcounter{page}{2}
\fancyfoot[c]{\thepage}
\fancyhead[L] {Работа № *labnum*} % some information in page header
\fancyhead[R]{}
	
	\section*{Ход работы и обработка данных}
	
	\subsection*{Определение разрешённых направлений поляроидов}
	
	\begin{enumerate}
		\item Разместим на оптической скамье осветитель
		$S$, поляроид $P_1$ и чёрное зеркало. Поворачивая поляроид вокруг направления луча, добьёмся наименьшей яркости отражённого пятна. Оставим поляроид в этом положении и вращением зеркала вокруг вертикальной оси снова добьёмся минимальной интенсивности отражённого луча.
		
		Для первого поляроида разрешённое направление горизонтальное, на лимбе $\alpha_1 \approx 78^{\circ} ~ \pm ~ 2^{\circ}$.
		
		\item Вместо чёрного зеркала поставим второй поляроид. Скрестим их, определим разрешённое направление второго поляроида -- горизонтальная волна, на лимбе $\alpha_2 \approx 36^{\circ} ~ \pm ~ 2^{\circ}$.
	\end{enumerate}

	\subsection*{Определение угла Брюстера для эбонита}
	\begin{enumerate}
		\item Поставим на скамью вместо чёрного зеркала эбонитовую пластину с круговой шкалой.
		
		\item  Повернем эбонитовое зеркало вокруг вертикальной оси так, чтобы его
		плоскость была перпендикулярна лучу, начальная метка отсчета на шкале $\alpha_0 = 214^{\circ}$.
		
		\item Установим направление разрешённых колебаний поляроида $P_1$ горизонтально и найдите угол поворота эбонита $\varphi_{\text{Б}}$, при котором интенсивность отражённого луча минимальна: его абсолютное значение равно $\varphi_{\text{Б}} = 273^{\circ} - 214^{\circ} = 59^{\circ}$.
		
		\item Повторим измерения, добавив светофильтр. Угол Брюстера в этом эксперименте оказался равным $\varphi_{\text{Б}} = 56^{\circ}$.
		
		\item По углу Брюстера рассчитаем показатель преломления.
		\begin{equation}
			n = \tan \varphi_{\text{Б}} \Rightarrow n = 1.66
		\end{equation}
		Табличное значение показателя преломления эбонита $n = 1.6$.
		
	\end{enumerate}

	\subsection*{Исследование стопы}
	
	Поставим стопу стеклянных пластинок вместо эбонитового зеркала и подберем для неё такое положение, при котором свет падает на стопу под углом Брюстера. Осветим стопу неполяризованным светом и, рассматривая через поляроиды свет, отражённый от стопы, определим ориентацию вектора
	\textbf{E} в отражённом луче и в преломлённом луче.
		
	\begin{center}
		Для отраженного света $\gamma_1 = 85^{\circ}$ \\
		Для прошедшего света $\gamma_2 = 171^{\circ}$
	\end{center} 

	\subsection*{Определение главных плоскостей двоякопреломляющих пластин}
	
	Поставим кристаллическую пластинку между скрещенными поляроидами. Вращая пластинку вокруг направления луча и наблюдая за интенсивностью света, про ходящего сквозь второй поляроид, определим, при каком условии главные направления пластинки совпадают с разрешёнными направлениями поляроидов. Повторим опыт для второй пластинки.
	
	\begin{center}
		Минимумы на первой пластинке наблюдаются при углах:\\
		$\alpha_1 = 12^{\circ} ~ \alpha_2 = 102^{\circ} ~ \alpha_3 = 190^{\circ} ~ \alpha_4 = 282^{\circ}$\\
		Минимумы на второй пластинке наблюдаются при углах:\\
		$\alpha_1 = 44^{\circ} ~ \alpha_2 = 134^{\circ} ~ \alpha_3 = 224^{\circ} ~ \alpha_4 = 314^{\circ}$\\
	\end{center}
	
	\subsection*{Выделение пластин $\lambda/2$ и $\lambda/4$}
	Добавим к схеме зелёный фильтр; установим разрешённое направление поляроида горизонтально, а главные направления исследуемой пластинки — под углом $45^{\circ}$ к горизонтали.
	
	Рассмотрим эффект от двух представленных пластинок:
	
	\begin{enumerate}
		\item Изучении первой пластинки: при вращении второго поляроида слегка (но не особо сильно) меняется цвет, а интенсивность света остается неизменной, то есть мы наблюдаем эллиптическую (или круговую поляризацию).
		\item Изучение второй пластинки: при вращении второго поляроида меняется интенсивность света, то есть мы наблюдаем линейную поляризацию.
	\end{enumerate} 

	Нам известно, что пластинка $\lambda/2$ не меняет характер поляризации, при её повороте, а пластинка $\lambda/4$ создаёт сдвиг фаз $\pi/2$ между колебаниями. Поскольку до прохождения пластинки свет был поляризован линейно (после прохождения через первый поляроид) изменение поляризации на эллиптическую (исчезновение минимумов при наблюдении) детектирует пластинку $\lambda/4$, сохранение минимумов (сохранение линейной поляризации) детектирует пластинку $\lambda/2$.
	
	\noindent \textbf{Вывод:} наблюдаемая пластинка №1 -- $\lambda/4$, пластинка №2 -- $\lambda/2$.
	
	\subsection*{Определение направлений большей и меньшей скоростей в пластинке $\lambda/4$}
	
	\begin{enumerate}
		\item Поставим между скрещенными поляроидами пластинку чувствительного оттенка ($\lambda$ для зелёного света), имеющую вид стрелки. Световой вектор, ориентированный вдоль направления стрелки, проходит с большей скоростью, перпендикулярный -- с меньшей.
		
		Установим разрешённое направление первого поляроида горизонтально и убедимся с помощью второго поляроида, что эта пластинка
		не меняет поляризацию зелёного света в условиях предыдущего опыта.
		
		\item Уберем зелёный фильтр и поставим между скрещенными поляроидами
		пластинку $\lambda$ (стрелка под углом $45^{\circ}$ к разрешённым направлениям поляроидов). Глядя сквозь второй поляроид на стрелку, убедимся, что она имеет пурпурный цвет (зелёный свет задерживается вторым поляроидом, а красная и синяя компоненты проходят).
		
		\item . Добавим к схеме пластинку $\lambda/4$, главные направления
		которой совпадают с главными направлениями пластины $\lambda$ и ориентированы под углом $45^{\circ}$ к разрешённым направлениям скрещенных поляроидов
		
		\item Теперь уберём пластину чувствительного оттенка. После второго поляроида интенсивность минимальная -- значит, быстрая ось пластинки направлена горизонтально, направление вращения правое, направление колебаний в первом и третьем квадрантах (разность фаз $\pi/4$). При повороте рейтера со стрелкой на $180^{\circ}$ вокруг вертикальной оси
		цвет стрелки меняется от зелёно-голубого до оранжево-жёлтого.
		
	\end{enumerate}
	
	\subsection*{Определение направления вращения светового вектора в эллиптически поляризованной волне}
	
	\begin{enumerate}
		\item  Снова поставим зелёный фильтр, а за ним между скрещенными поляроидами -- пластинку произвольной толщины ($\lambda/4$).
		
		\item Получим эллиптически-поляризованный свет. Для этого установим разрешённое направление первого поляроида под углом $10-20^{\circ}$ к горизонтали так, чтобы вектор \textbf{E} падающего на пластинку света был расположен в первом квадранте.
		Установим разрешённое направление второго поляроида вертикально и, вращая пластинку, найдем минимальную
		интенсивность света, прошедшего второй поляроид. Вращая второй поляроид, убедитесь, что свет поляризован эллиптически,
		а не линейно.
		Таким образом, получим эллипс поляризации с вертикально ориентированной малой осью.
		\item  Для определения направления вращения светового вектора в эллипсе
		установим между поляроидами дополнительную пластинку $\lambda/4$ с известными направлениями <<быстрой>> и <<медленной>> осей, ориентированными по осям эллипса поляризации анализируемого света.
		В этом случае вектор \textbf{E} на выходе будет таким, как если бы свет прошёл две пластинки $\lambda/4$: свет на выходе из второй пластинки будет линейно поляризован. Если пластинки поодиночке дают эллипсы, вращающиеся в разные стороны, то поставленные друг за другом, они скомпенсируют
		разность фаз, и вектор \textbf{E} на выходе останется в первом
		и третьем квадрантах. Если же световой вектор перешёл в смежные квадранты, значит, эллипсы вращаются в одну сторону. 
		
		После второго поляроида интенсивность света максимальна. Значит, две пластины усиливают друг друга, световой вектор перешёл в смежные квадранты, эллипсы вращаются в одну сторону.
		
	\end{enumerate}
	
	\subsection*{Интерференция поляризованных лучей}
	
	Расположим между скрещенными поляроидами мозаичную слюдяную пластинку. Она собрана из 4-х узких полосок слюды, лежащих по сторонам квадрата (две полоски $\lambda/4$ и по одной -- $\lambda/2$ и $3\lambda/4$).
	 
	Вращаем пластинку: изменяется интенсивность света с периодичностью $\pi/4$
	
	Вращаем второй поляроид: изменяется цвет пластинок также с периодичностью  $\pi/4$.

\end{document}